\documentclass[11pt,a4paper]{article}
\usepackage[margin=1in]{geometry}
\usepackage{times}
\usepackage{latexsym}
\usepackage{graphicx}
\usepackage{amsmath}
\usepackage{amssymb}
\usepackage{booktabs}
\usepackage{multirow}
\usepackage{xcolor}
\usepackage{url}
\usepackage{hyperref}
\usepackage{tabularx}
\usepackage{microtype}
\usepackage{subcaption}
\usepackage{natbib}

\title{When Data Programming Meets Agent Trajectories:\\An Empirical Study of Weak Supervision Failure Modes\\for Progress Stagnation Detection}

\author{Anonymous}

\date{}

\begin{document}
\maketitle

%==============================================================================
% ABSTRACT
%==============================================================================
\begin{abstract}
LLM-based agents frequently enter \emph{progress stagnation}---executing surface-diverse but collectively ineffective action sequences during long-horizon tasks.
Detecting such stagnation is important for cost control and intervention, yet the concept itself resists clean operationalization: human annotators achieve only moderate agreement ($\kappa{=}0.52$).
We formulate stagnation detection as a \emph{window-level} binary classification problem and build a weak-supervision pipeline with four labeling functions (LFs) and a GRU-based probe (ProgressProbe) trained on the resulting noisy labels.
Rather than claiming methodological advances, we present a \emph{systematic empirical study} that exposes three categories of failure modes when applying data programming to sequential agent data:
(1)~LFs exhibit systematic temporal dynamics (leading/lagging indicators) that violate the i.i.d.\ assumption of standard data programming, yet a temporal label model fails to outperform FlyingSquid ($\Delta{=}{-}0.2$pp on the gold set);
(2)~the strongest LF (\texttt{outcome\_backpropagation}) leaks trajectory-level outcome information, inflating label quality by 7.4pp;
(3)~cross-agent transfer degrades severely (F1 drops from 0.91 to 0.61).
On the positive side, GRU temporal modeling shows robust performance on noisy labels ($+$6.9pp over mean-pooling at conf${\geq}$1), and provides a 615$\times$ cheaper alternative to LLM-as-judge for continuous monitoring (gold-set F1 0.72 vs.\ 0.86 for Claude CoT).
Our stagnation taxonomy, window-level formulation, and failure-mode characterization provide empirical foundations for future work on agent progress monitoring.
\end{abstract}

%==============================================================================
% 1. INTRODUCTION
%==============================================================================
\section{Introduction}

LLM-based autonomous agents~\cite{yang2024swe,wang2024opendevin} increasingly tackle long-horizon tasks such as software engineering, where a single trajectory may span dozens of tool-use steps.
A recurring failure mode is \emph{progress stagnation}: the agent enters a regime of surface-diverse but goal-unrelated actions---trying different commands, editing files, reading documentation---without making meaningful progress toward the objective.
Unlike simple repetition loops, stagnation is semantically subtle: each individual action may appear reasonable, and only the \emph{window-level} pattern reveals the lack of forward momentum.

Detecting stagnation in real time would enable early termination of futile trajectories, saving computational resources and enabling meta-controllers to intervene.
However, three fundamental challenges impede reliable detection:
(a)~stagnation is inherently ambiguous at the boundary between exploration and futility (our human annotators achieve only $\kappa{=}0.52$);
(b)~agent trajectories violate the i.i.d.\ assumption underlying standard data programming~\cite{ratner2017snorkel}, as labeling functions exhibit systematic temporal dependencies;
(c)~trajectory-level outcome information can leak into window-level labels, creating inflated performance estimates.

We formulate stagnation detection as window-level binary classification over $k$-step windows of agent trajectories, and construct a weak-supervision pipeline consisting of four labeling functions and a GRU-based neural probe.
Our original hypothesis was that a temporal label model exploiting LF dynamics would yield better labels than standard aggregation, and that the trained probe would exceed the LF ensemble (``student exceeds teacher'').
\textbf{Neither hypothesis holds.}
Instead, we find that:

\begin{enumerate}
\item A temporal label model (HMM-based) achieves 0.770 gold-set F1, statistically indistinguishable from FlyingSquid's 0.772, despite successfully recovering leading/lagging LF structure. Correlation-aware aggregation already captures the relevant dependencies without explicit temporal modeling.
\item The probe achieves 0.72 gold-set F1, \emph{below} the LF oracle (0.745) rather than above it. High auto-label performance (0.91) partially reflects overfitting to LF noise patterns rather than learning genuine progress signals.
\item LLM-as-judge with chain-of-thought reasoning achieves 0.855 gold-set F1, clearly outperforming the probe by 13.5pp. The probe's advantage is limited to cost (615$\times$ cheaper) and latency (667$\times$ faster).
\item Cross-agent transfer from SWE-agent to OpenHands degrades probe F1 from 0.91 to 0.61 ($-$29.4pp).
\end{enumerate}

We reframe these negative results as contributions to the empirical understanding of weak supervision on sequential agent data:
a stagnation taxonomy and window-level formulation as a task contribution (\S\ref{sec:task});
systematic characterization of data programming failure modes (\S\ref{sec:failure-modes});
and evidence that GRU temporal modeling provides robustness to label noise (\S\ref{sec:robustness}).

%==============================================================================
% 2. RELATED WORK
%==============================================================================
\section{Related Work}

\paragraph{Agent monitoring and intervention.}
Cycle detectors such as CDCS+CDSA~\cite{cycles2024} and SpecRA~\cite{specra2024} identify structural repetition but miss semantically diverse stagnation.
Self-reflection approaches~\cite{shinn2023reflexion,yao2024tree} rely on the trapped model's own judgment.
Process reward models~\cite{lightman2023lets,wang2024math} score individual steps but require step-level outcome labels and lose window-level temporal dynamics.
Value-function methods~\cite{le2022coderl} need shaped rewards.
Our work differs by framing progress monitoring as a weak-supervision problem at the window level.

\paragraph{Data programming and weak supervision.}
Snorkel~\cite{ratner2017snorkel} and its successors~\cite{fu2020fast,shin2022universalizing} assume conditionally independent labeling functions.
FlyingSquid~\cite{fu2020fast} relaxes this to model pairwise correlations via the triplet method.
These approaches assume i.i.d.\ data points; we empirically characterize how this assumption fails on sequential agent trajectories with temporally structured LF outputs.

\paragraph{LLM-as-judge.}
Using LLMs as evaluators~\cite{zheng2024judging} provides high-quality assessments but at substantial cost.
We compare against LLM-as-judge baselines to establish an upper bound on detection quality.

%==============================================================================
% 3. TASK FORMULATION
%==============================================================================
\section{Task Formulation}
\label{sec:task}

\subsection{Stagnation Taxonomy}

We define \emph{progress stagnation} as a contiguous window of agent actions that fails to advance the agent toward its declared objective, despite surface-level activity.
Through analysis of 4,880 SWE-agent trajectories from SWE-bench Verified, we identify three stagnation subtypes: (1)~\emph{repetitive stagnation}---the agent repeats structurally similar commands; (2)~\emph{exploratory stagnation}---the agent tries diverse approaches without convergence; (3)~\emph{partial-progress stagnation}---some file edits occur but produce no functional improvement.

\subsection{Window-Level Formulation}

Given a trajectory $\mathcal{T} = (a_1, a_2, \ldots, a_T)$ of agent actions, we extract overlapping windows $w_i = (a_i, a_{i+1}, \ldots, a_{i+k-1})$ of size $k$.
Each window is labeled as \textsc{Progress} or \textsc{Stagnation}.
We set $k{=}5$ based on ablation (Table~\ref{tab:window-ablation}), yielding 54,809 windows from 4,880 trajectories.

\subsection{Labeling Functions}
\label{sec:lfs}

We design four labeling functions that capture complementary signals:

\begin{itemize}
\item \textbf{command\_repetition}: Flags windows where $\geq$3 of 5 steps execute structurally similar commands. Leading indicator (fires ${\sim}$1.8 steps before stagnation onset).
\item \textbf{output\_similarity}: Measures pairwise cosine similarity of step outputs within the window. Concurrent indicator.
\item \textbf{edit\_cycle\_detection}: Detects file-edit cycles (edit $\to$ revert $\to$ re-edit patterns). Concurrent indicator.
\item \textbf{outcome\_backpropagation}: Propagates trajectory-level outcome (resolved/unresolved) backward to label late windows. Lagging indicator (fires ${\sim}$2.1 steps after onset).
\end{itemize}

Labels are aggregated via majority vote.
We define label \emph{confidence} as the number of agreeing LFs: conf${\geq}$2 yields 10,780 high-confidence windows; conf${\geq}$1 yields 38,161 windows (72\% at conf${=}$1 with substantial LF disagreement).

\subsection{Human Annotation (Gold Set)}
\label{sec:gold}

Three annotators labeled 420 windows (4.2 min/window average; 29.4 hours total).
Table~\ref{tab:human-annotation} summarizes inter-annotator agreement.

\begin{table}[t]
\centering
\small
\begin{tabular}{lr}
\toprule
\textbf{Metric} & \textbf{Value} \\
\midrule
Windows annotated & 420 \\
Fleiss $\kappa$ & 0.52 \\
Unanimous agreement & 45\% \\
Majority agreement & 68\% \\
Auto MV vs.\ gold agreement & 73\% \\
Gold stagnation prevalence & 38\% \\
\bottomrule
\end{tabular}
\caption{Human annotation statistics. Moderate $\kappa$ reflects inherent ambiguity of the progress/stagnation boundary.}
\label{tab:human-annotation}
\end{table}

Disagreement analysis (Table~\ref{tab:disagreement}) reveals that 43\% of annotator disagreements involve partial-progress cases, and 26\% concern exploration-vs-stagnation distinctions. Notably, 16\% of disagreements are influenced by annotator knowledge of the trajectory's final outcome, mirroring the \texttt{outcome\_bp} leakage issue in our LFs.

\begin{table}[t]
\centering
\small
\begin{tabular}{lrl}
\toprule
\textbf{Type} & \textbf{\%} & \textbf{Pattern} \\
\midrule
Partial progress & 43\% & Edits but no test improvement \\
Exploration vs stag. & 26\% & Trying different approaches \\
Outcome-influenced & 16\% & Knowledge of final result \\
Annotation error & 15\% & Missed context / fatigue \\
\bottomrule
\end{tabular}
\caption{Breakdown of annotator disagreements ($n{=}135$ disagreement cases out of 420 windows).}
\label{tab:disagreement}
\end{table}

%==============================================================================
% 4. METHODS
%==============================================================================
\section{Methods}

\subsection{Label Models}
\label{sec:label-models}

\paragraph{Majority Vote (MV).} Standard unweighted vote across four LFs.

\paragraph{FlyingSquid.} Correlation-aware label model~\cite{fu2020fast} using the triplet method to estimate pairwise LF dependencies without temporal structure.

\paragraph{Temporal Label Model.} We propose an HMM-based label model that explicitly incorporates LF temporal roles.
The model assigns each LF a temporal weight based on its lead/lag relative to stagnation onset, learned via EM on the LF output sequences.
Two variants: (1)~HMM with 2 hidden states; (2)~weighted vote with learned temporal weights.

\subsection{ProgressProbe}

ProgressProbe is a lightweight GRU-based binary classifier.
Each step's action text is encoded via Qwen3-Embedding-0.6B (1024d).
A single-layer GRU (hidden size 128) processes the $k$-step embedding sequence, and a linear head classifies the final hidden state.
Training uses cross-entropy loss on weak labels; we also evaluate a SupCon auxiliary loss variant.
Total parameters: 198K (GRU + mean-pool) or 285K (GRU + attention).

\subsection{Baselines}

\begin{itemize}
\item \textbf{Repeat-3}: Flags stagnation if $\geq$3 consecutive steps share $>$90\% command overlap.
\item \textbf{CosineSim}: Thresholds mean pairwise cosine similarity within window (optimized threshold $=$ 0.78).
\item \textbf{Mean-Pool + LR/MLP}: Logistic regression or 2-layer MLP on mean-pooled step embeddings (no temporal modeling).
\item \textbf{Structural (GBT)}: Gradient-boosted trees on 6 non-NLP features (lines changed, file count, test count, command-type entropy, AST diff size, edit frequency).
\item \textbf{LF Oracle}: Direct application of LF ensemble at inference time (majority vote on real-time LF outputs).
\item \textbf{LLM-as-Judge}: Claude 3.5 Sonnet and GPT-4o in zero-shot, few-shot (3 examples), and chain-of-thought (CoT) settings.
\end{itemize}

%==============================================================================
% 5. EXPERIMENTS
%==============================================================================
\section{Experiments}

\subsection{Main Results on Gold Set}

Table~\ref{tab:main-results} presents all methods evaluated on the 420-window human-annotated gold set.

\begin{table}[t]
\centering
\small
\begin{tabular}{lccc}
\toprule
\textbf{Method} & \textbf{Gold F1} & \textbf{Cost/win.} & \textbf{Latency} \\
\midrule
\multicolumn{4}{l}{\emph{LLM-as-Judge}} \\
\quad Claude CoT & 0.855 & \$0.042 & 3.8s \\
\quad Claude few-shot & 0.820 & \$0.035 & 2.9s \\
\quad GPT-4o CoT & 0.840 & \$0.038 & 3.5s \\
\quad GPT-4o few-shot & 0.808 & \$0.032 & 2.7s \\
\midrule
\multicolumn{4}{l}{\emph{Label Model / Oracle}} \\
\quad LF Oracle (MV) & 0.745 & $\sim$0 & $\sim$0 \\
\quad FlyingSquid & 0.772 & $\sim$0 & $\sim$0 \\
\quad Temporal LM (HMM) & 0.770 & $\sim$0 & $\sim$0 \\
\midrule
\multicolumn{4}{l}{\emph{Neural Probes}} \\
\quad ProgressProbe (GRU) & 0.720 & \$5.2e-5 & 4.2ms \\
\quad Mean-Pool MLP & 0.701 & \$5.2e-5 & 1.8ms \\
\quad Mean-Pool LR & 0.688 & \$5.2e-5 & 0.5ms \\
\midrule
\multicolumn{4}{l}{\emph{Heuristic / Structural}} \\
\quad Structural (GBT) & 0.680 & $\sim$0 & 0.3ms \\
\quad Repeat-3 & --- & $\sim$0 & $\sim$0 \\
\quad CosineSim (opt.) & 0.582 & $\sim$0 & 0.8ms \\
\bottomrule
\end{tabular}
\caption{Main results on gold set. LLM-as-judge methods set the quality upper bound; the probe offers a 615$\times$ cost advantage at reduced accuracy. The probe does not exceed the LF oracle.}
\label{tab:main-results}
\end{table}

Three findings stand out.
First, LLM-as-judge with CoT clearly dominates in accuracy (Claude CoT: 0.855), establishing an upper bound that the probe (0.720) falls 13.5pp short of.
Second, the probe does not exceed the LF oracle (0.745 vs.\ 0.720, $\Delta{=}{-}2.5$pp), invalidating the ``student exceeds teacher'' hypothesis.
Third, the probe's practical advantage is limited to cost (615$\times$ cheaper than LLM-as-judge) and latency (667$\times$ faster), making it suitable only for high-volume continuous monitoring where LLM costs are prohibitive.

\subsection{Multi-Seed Stability}

Table~\ref{tab:multiseed} shows probe performance across three random seeds, revealing the gap between auto-label and gold-set evaluation.

\begin{table}[t]
\centering
\small
\begin{tabular}{lcccc}
\toprule
\textbf{Seed} & \textbf{Auto F1} & \textbf{Gold F1} & \textbf{$\Delta$ Repeat-3} & \textbf{$\Delta$ LF Oracle} \\
\midrule
42 & 0.906 & 0.738 & +6.2pp & $-$0.8pp \\
123 & 0.887 & 0.692 & +4.4pp & $-$2.5pp \\
456 & 0.915 & 0.725 & +7.2pp & $-$1.2pp \\
\midrule
Mean$\pm$Std & 0.903$\pm$.019 & 0.718$\pm$.031 & +5.9$\pm$1.4 & $-$1.5$\pm$0.9 \\
\bottomrule
\end{tabular}
\caption{Multi-seed stability. Auto-label F1 is stable (std 0.019), but gold-set F1 shows greater variance (std 0.031) and consistently falls below the LF oracle.}
\label{tab:multiseed}
\end{table}

The divergence between auto-label performance (0.903) and gold-set performance (0.718) indicates that part of the probe's auto-label success reflects overfitting to LF noise patterns rather than learning true progress signals.

\subsection{Label Model Comparison}
\label{sec:label-model-comparison}

Table~\ref{tab:label-model} compares label aggregation strategies on the gold set. The temporal label model (HMM) achieves 0.770 gold-set F1, indistinguishable from FlyingSquid (0.772, $\Delta{=}{-}0.2$pp). Both improve over majority vote by ${\sim}$2.5pp.

\begin{table}[t]
\centering
\small
\begin{tabular}{lcccc}
\toprule
\textbf{Label Model} & \textbf{Gold F1} & \textbf{Stag F1} & \textbf{Prog F1} & \textbf{Flip Rate} \\
\midrule
Majority Vote & 0.745 & 0.680 & 0.810 & --- \\
FlyingSquid & 0.772 & 0.712 & 0.832 & 11.2\% \\
Temporal LM (HMM) & 0.770 & 0.708 & 0.832 & 8.3\% \\
Temporal LM (weighted) & 0.762 & 0.698 & 0.826 & 6.5\% \\
\bottomrule
\end{tabular}
\caption{Label model comparison. Temporal label models do not improve over FlyingSquid despite recovering temporal structure.}
\label{tab:label-model}
\end{table}

The HMM successfully learns leading/lagging weights (Table~\ref{tab:temporal-weights}), but these adjustments flip only 8.3\% of labels, with flipped-label accuracy on the gold set of approximately 55\%---near random. FlyingSquid's correlation-aware aggregation already captures the LF dependencies without explicit temporal modeling, because the temporal structure manifests as statistical correlations that the triplet method recovers.

\begin{table}[t]
\centering
\small
\begin{tabular}{lccc}
\toprule
\textbf{LF} & \textbf{Role} & \textbf{Lag (steps)} & \textbf{Weight} \\
\midrule
command\_repetition & Leading & $-$1.8 & 1.35 \\
output\_similarity & Concurrent & 0.0 & 1.12 \\
edit\_cycle\_detection & Concurrent & $-$0.3 & 1.08 \\
outcome\_backprop. & Lagging & +2.1 & 0.72 \\
\bottomrule
\end{tabular}
\caption{Learned temporal weights from the HMM label model. \texttt{command\_repetition} is a leading indicator; \texttt{outcome\_bp} is lagging.}
\label{tab:temporal-weights}
\end{table}

%==============================================================================
% 5.4 LF TEMPORAL DYNAMICS
%==============================================================================
\subsection{LF Temporal Dynamics}
\label{sec:lf-dynamics}

Table~\ref{tab:lf-temporal} characterizes the temporal roles of each LF, confirming that LFs on agent trajectories violate the i.i.d.\ assumption of standard data programming.

\begin{table}[t]
\centering
\small
\begin{tabular}{lcccc}
\toprule
\textbf{LF} & \textbf{Lead/Lag} & \textbf{Role} & \textbf{Autocorr.} & \textbf{Onset Rate} \\
\midrule
cmd\_repetition & $-$1.8 & Leading & 0.72 & 0.68 \\
output\_similarity & $-$0.5 & Concurrent & 0.65 & 0.61 \\
edit\_cycle\_det. & $-$0.2 & Concurrent & 0.58 & 0.55 \\
outcome\_bp & +2.1 & Lagging & 0.85 & 0.78 \\
\bottomrule
\end{tabular}
\caption{LF temporal dynamics. High autocorrelation ($>$0.58) indicates slow-varying LF outputs; the leading/lagging structure is systematic but diluted within $k{=}5$ windows.}
\label{tab:lf-temporal}
\end{table}

\texttt{command\_repetition} fires ${\sim}$1.8 steps before stagnation onset, while \texttt{outcome\_bp} fires ${\sim}$2.1 steps after.
However, high lag-1 autocorrelation (0.72--0.85) means LF outputs change slowly, and within a 5-step window, leading/lagging differences are averaged out.
LF transition consistency is only 0.61---insufficient for reliable temporal label adjustment.
This explains why the temporal label model gains no advantage over FlyingSquid: the temporal signal exists but is too weak at the window granularity to improve label aggregation.

%==============================================================================
% 5.5 FAILURE MODE: OUTCOME LEAKAGE
%==============================================================================
\subsection{Failure Mode: Outcome Leakage}
\label{sec:failure-modes}

\paragraph{Outcome backpropagation leakage.}
The \texttt{outcome\_bp} LF propagates trajectory-level outcome information backward to window-level labels.
Table~\ref{tab:outcome-bp} quantifies this leakage.

\begin{table}[t]
\centering
\small
\begin{tabular}{lcccc}
\toprule
\textbf{Variant} & \textbf{Macro-F1} & \textbf{Stag F1} & \textbf{Prog F1} & \textbf{$\Delta$} \\
\midrule
Original (all LFs) & 0.906 & 0.859 & 0.952 & --- \\
Hold-out outcome\_bp & 0.832 & 0.762 & 0.902 & $-$7.4pp \\
outcome\_bp last-3 only & 0.868 & 0.815 & 0.921 & $-$3.8pp \\
outcome\_bp last-1 only & 0.812 & 0.738 & 0.886 & $-$9.4pp \\
\bottomrule
\end{tabular}
\caption{Outcome backpropagation sensitivity. Holding out \texttt{outcome\_bp} drops auto-label F1 by 7.4pp, with stagnation F1 most affected ($-$9.7pp).}
\label{tab:outcome-bp}
\end{table}

Exit context is present in 78\% of trajectories, affecting 31\% of windows.
The \texttt{outcome\_bp} marginal contribution (0.074) is the largest among all LFs (Table~\ref{tab:lf-holdout}), yet it encodes ``whether the task ultimately succeeded'' rather than ``whether this window makes progress.''
This leakage inflates the entire pipeline's performance: the ``clean'' performance (without \texttt{outcome\_bp}) is 0.832 on auto-labels, and likely 0.68--0.70 on the gold set.

\paragraph{LF hold-one-out analysis.}
Table~\ref{tab:lf-holdout} shows marginal contributions of each LF.
\texttt{outcome\_bp} dominates the ensemble, contributing 7.4pp---more than double the next most important LF.

\begin{table}[t]
\centering
\small
\begin{tabular}{lccc}
\toprule
\textbf{Held-out LF} & \textbf{Auto F1} & \textbf{$\Delta$} & \textbf{Marginal} \\
\midrule
None (full) & 0.906 & --- & --- \\
command\_repetition & 0.875 & $-$3.1pp & 0.031 \\
output\_similarity & 0.882 & $-$2.4pp & 0.024 \\
edit\_cycle\_detection & 0.891 & $-$1.5pp & 0.015 \\
outcome\_backpropagation & 0.832 & $-$7.4pp & 0.074 \\
\bottomrule
\end{tabular}
\caption{LF hold-one-out. The leakiest LF (\texttt{outcome\_bp}) is also the most influential---a cautionary pattern for weak supervision practitioners.}
\label{tab:lf-holdout}
\end{table}

%==============================================================================
% 5.6 CROSS-AGENT TRANSFER
%==============================================================================
\subsection{Failure Mode: Cross-Agent Transfer}
\label{sec:cross-agent}

We evaluate all methods on 2,340 OpenHands trajectories (18,720 windows) to test cross-implementation transfer from SWE-agent.
Table~\ref{tab:cross-agent} reveals severe degradation across all methods.

\begin{table}[t]
\centering
\small
\begin{tabular}{lccc}
\toprule
\textbf{Method} & \textbf{Same-Agent} & \textbf{Cross-Agent} & \textbf{Degrad.} \\
\midrule
Probe (GRU) & 0.906 & 0.612 & $-$29.4pp \\
LF Oracle (MV) & 0.945 & 0.685 & $-$26.0pp \\
Repeat-3 & 0.843 & 0.590 & $-$25.3pp \\
Mean-Pool LR & 0.843 & 0.558 & $-$28.5pp \\
CosineSim & 0.582 & 0.480 & $-$10.2pp \\
\bottomrule
\end{tabular}
\caption{Cross-agent transfer results. The probe degrades more than the LF oracle ($-$29.4 vs.\ $-$26.0pp), contradicting the hypothesis that it learns transferable progress signals.}
\label{tab:cross-agent}
\end{table}

The probe's degradation ($-$29.4pp) exceeds even the LF oracle's ($-$26.0pp), directly contradicting the hypothesis that the probe learns agent-agnostic progress signals.
LF coverage on OpenHands drops to 82\% (vs.\ 95\% on SWE-agent), with \texttt{outcome\_bp} coverage declining most sharply ($-$16\%).
The LF oracle---despite also degrading---outperforms the probe on cross-agent evaluation by 7.3pp, suggesting that heuristic rules transfer more robustly than learned features.

\begin{table}[t]
\centering
\small
\begin{tabular}{lccc}
\toprule
\textbf{LF} & \textbf{SWE-agent} & \textbf{OpenHands} & $\Delta$ \\
\midrule
command\_repetition & 0.95 & 0.88 & $-$7\% \\
output\_similarity & 0.92 & 0.85 & $-$7\% \\
outcome\_bp & 0.78 & 0.62 & $-$16\% \\
edit\_cycle\_detection & 0.88 & 0.79 & $-$9\% \\
Ensemble & 0.95 & 0.82 & $-$13\% \\
\bottomrule
\end{tabular}
\caption{LF coverage across agent implementations. \texttt{outcome\_bp} suffers the largest drop.}
\label{tab:lf-coverage}
\end{table}

%==============================================================================
% 5.7 LLM-AS-JUDGE COMPARISON
%==============================================================================
\subsection{LLM-as-Judge}

Table~\ref{tab:llm-judge} provides the full comparison. LLM-as-judge with CoT reasoning achieves substantially higher gold-set F1 than the probe, with Claude CoT reaching 0.855 ($+$13.5pp over probe).

\begin{table}[t]
\centering
\small
\begin{tabular}{llcccc}
\toprule
\textbf{Model} & \textbf{Setting} & \textbf{Gold F1} & \textbf{Stag} & \textbf{Prog} & \textbf{Cost} \\
\midrule
Claude 3.5 & zero-shot & 0.780 & 0.725 & 0.835 & \$0.028 \\
Claude 3.5 & few-shot & 0.820 & 0.772 & 0.868 & \$0.035 \\
Claude 3.5 & CoT & 0.855 & 0.810 & 0.900 & \$0.042 \\
GPT-4o & zero-shot & 0.765 & 0.710 & 0.820 & \$0.025 \\
GPT-4o & few-shot & 0.808 & 0.758 & 0.858 & \$0.032 \\
GPT-4o & CoT & 0.840 & 0.795 & 0.885 & \$0.038 \\
\midrule
Probe & --- & 0.720 & 0.650 & 0.790 & \$5.2e-5 \\
\bottomrule
\end{tabular}
\caption{LLM-as-judge vs.\ probe on gold set. LLMs achieve higher accuracy via semantic understanding of goal-directed progress; the probe's advantage is purely economic.}
\label{tab:llm-judge}
\end{table}

At full scale (54K windows), LLM CoT evaluation costs ${\sim}$\$1,750 versus ${\sim}$\$2.80 for the probe.
For continuous monitoring of millions of windows, LLM evaluation is economically infeasible, making the probe the only viable option despite its lower accuracy.

%==============================================================================
% 5.8 ERROR ANALYSIS
%==============================================================================
\subsection{Error Analysis}

We analyze 112 probe errors on the gold set (Table~\ref{tab:error-analysis}).

\begin{table}[t]
\centering
\small
\begin{tabular}{lccl}
\toprule
\textbf{Error Type} & \textbf{N} & \textbf{\%} & \textbf{Pattern} \\
\midrule
FN: Semantic-diverse stag. & 37 & 33 & Surface-diverse, goal-unrelated \\
FN: Partial progress & 17 & 15 & Mixed progress and stagnation \\
FN: Other & 10 & 9 & Edge cases \\
FP: Exploratory actions & 25 & 22 & Trial-and-error $\approx$ stagnation \\
FP: conf=1 region & 35 & 31 & LF disagreement zone \\
\bottomrule
\end{tabular}
\caption{Error breakdown. The dominant failure mode (33\%) is semantically diverse stagnation, which requires goal-level reasoning beyond embedding statistics.}
\label{tab:error-analysis}
\end{table}

False negatives dominate (64 vs.\ 48 FPs).
The largest FN category (33\%) involves semantically diverse stagnation: the agent executes varied but collectively futile actions.
These cases require goal-level reasoning to detect---exactly where LLM CoT excels and embedding-based probes fail.
Among FPs, 73\% fall in conf${=}$1 regions (high LF disagreement), and 35\% depend on \texttt{outcome\_bp}, confirming that FPs are concentrated in the noisiest parts of the label space.

%==============================================================================
% 5.9 ROBUSTNESS TO LABEL NOISE
%==============================================================================
\subsection{GRU Robustness to Label Noise}
\label{sec:robustness}

Table~\ref{tab:noise-robustness} compares training on high-confidence (conf${\geq}$2) versus full (conf${\geq}$1) data.

\begin{table}[t]
\centering
\small
\begin{tabular}{lcccc}
\toprule
\textbf{Training Set} & \textbf{N} & \textbf{Auto F1} & \textbf{Gold F1} & \textbf{Stag F1} \\
\midrule
conf$\geq$2 & 10,780 & 0.906 & 0.720 & 0.859 \\
conf$\geq$1 (full) & 38,161 & 0.847 & 0.735 & 0.777 \\
$\Delta$ & +27,381 & $-$5.9pp & +1.5pp & $-$8.2pp \\
\bottomrule
\end{tabular}
\caption{Label noise tolerance. The GRU trained on noisy full data (conf$\geq$1) achieves \emph{higher} gold-set F1 than on clean conf$\geq$2 data, despite lower auto-label F1.}
\label{tab:noise-robustness}
\end{table}

A notable reversal occurs: while auto-label F1 drops 5.9pp on noisy data (expected), gold-set F1 \emph{increases} by 1.5pp.
The larger training set compensates for noise, and the GRU's temporal modeling helps extract signal from noisy labels.
This effect is amplified when comparing GRU vs.\ mean-pooling on noisy data: GRU gains +6.87pp over mean-pool MLP at conf${\geq}$1 (Table~\ref{tab:temporal-vs-meanpool}), compared to +4.83pp at conf${\geq}$2, confirming that temporal modeling provides disproportionate benefit under high label noise.

\begin{table}[t]
\centering
\small
\begin{tabular}{lcccc}
\toprule
\textbf{Setting} & \textbf{Pool LR} & \textbf{Pool MLP} & \textbf{GRU} & \textbf{GRU gain} \\
\midrule
Auto conf$\geq$2 & 0.844 & 0.857 & 0.906 & +4.83pp \\
Auto conf$\geq$1 & 0.759 & 0.778 & 0.847 & +6.87pp \\
Gold set & 0.688 & 0.701 & 0.720 & +1.9pp \\
\bottomrule
\end{tabular}
\caption{GRU vs.\ mean-pooling. Temporal modeling's advantage grows under noisier labels (conf$\geq$1).}
\label{tab:temporal-vs-meanpool}
\end{table}

%==============================================================================
% 5.10 ABLATIONS
%==============================================================================
\subsection{Ablations}

\paragraph{Window size.}
Table~\ref{tab:window-ablation} shows $k{=}5$ is optimal, though gold-set differences are small (1.5pp range).

\begin{table}[t]
\centering
\small
\begin{tabular}{lccccc}
\toprule
$k$ & $N$ & \textbf{Auto F1} & \textbf{Stag F1} & \textbf{Prog F1} & \textbf{Gold F1} \\
\midrule
3 & 82,214 & 0.878 & 0.812 & 0.944 & 0.705 \\
5 & 54,809 & 0.906 & 0.859 & 0.952 & 0.720 \\
7 & 39,149 & 0.895 & 0.848 & 0.942 & 0.718 \\
\bottomrule
\end{tabular}
\caption{Window size ablation. $k{=}5$ balances signal accumulation and onset dilution.}
\label{tab:window-ablation}
\end{table}

\paragraph{Architecture.}
Table~\ref{tab:arch-ablation} compares architectures. Gold-set differences span only 2.7pp, suggesting label quality---not architecture---is the bottleneck.

\begin{table}[t]
\centering
\small
\begin{tabular}{lcccc}
\toprule
\textbf{Architecture} & \textbf{Params} & \textbf{Auto F1} & \textbf{Gold F1} & \textbf{ms} \\
\midrule
GRU + Attention & 285K & 0.912 & 0.728 & 5.1 \\
GRU + Mean-pool & 198K & 0.906 & 0.720 & 4.2 \\
Transformer (2L) & 1.24M & 0.903 & 0.722 & 8.7 \\
MLP on pooled & 132K & 0.857 & 0.701 & 1.8 \\
\bottomrule
\end{tabular}
\caption{Architecture ablation. Gold-set performance is insensitive to architecture choice (2.7pp range), confirming the bottleneck is label quality.}
\label{tab:arch-ablation}
\end{table}

\paragraph{Encoder.}
Four embedding encoders produce a 1.8pp range on the gold set (Table~\ref{tab:encoder-ablation}), confirming encoder choice is not the performance bottleneck.

\begin{table}[t]
\centering
\small
\begin{tabular}{lcccc}
\toprule
\textbf{Encoder} & \textbf{Dim} & \textbf{Auto F1} & \textbf{Gold F1} & \textbf{min} \\
\midrule
Qwen3-Emb-0.6B & 1024 & 0.906 & 0.720 & 53 \\
BGE-M3 & 1024 & 0.898 & 0.715 & 48 \\
E5-large-v2 & 1024 & 0.892 & 0.708 & 45 \\
GTE-base & 768 & 0.885 & 0.702 & 32 \\
\bottomrule
\end{tabular}
\caption{Encoder ablation.}
\label{tab:encoder-ablation}
\end{table}

\paragraph{SupCon auxiliary loss.}
Adding supervised contrastive loss yields +0.6pp on auto-labels and +0.3pp on the gold set---negligible.
Embedding separation improves (0.74$\to$0.82) but does not translate into meaningful classification gains when labels are noisy.

\paragraph{NL vs.\ structural features.}
NL embeddings outperform non-NLP structural features (GBT on 6 numeric features) by 4.0pp on the gold set (0.720 vs.\ 0.680), confirming that semantic content provides signal beyond structural patterns, though the gap is moderate.

\paragraph{Data scaling.}
Performance saturates at ${\sim}$50\% of training data: gold-set F1 reaches 0.715 at 50\% (5,400 windows) vs.\ 0.720 at 100\%, indicating the bottleneck is label quality, not quantity.

\begin{table}[t]
\centering
\small
\begin{tabular}{lccc}
\toprule
\textbf{Train \%} & \textbf{N} & \textbf{Auto F1} & \textbf{Gold F1} \\
\midrule
10\% & $\sim$1,080 & 0.862 & 0.695 \\
25\% & $\sim$2,700 & 0.882 & 0.708 \\
50\% & $\sim$5,400 & 0.898 & 0.715 \\
100\% & $\sim$10,780 & 0.906 & 0.720 \\
\bottomrule
\end{tabular}
\caption{Data scaling. Early saturation ($\sim$50\%) confirms label quality as the bottleneck.}
\label{tab:data-scaling}
\end{table}

\paragraph{CosineSim baseline.}
The original CosineSim baseline (F1$=$0.313) suffered from threshold miscalibration (default 0.5 in a high-similarity embedding space).
Threshold sweep identifies optimal threshold$=$0.78, yielding corrected F1$=$0.582.

\begin{table}[t]
\centering
\small
\begin{tabular}{lcccc}
\toprule
\textbf{Threshold} & \textbf{Prec.} & \textbf{Recall} & \textbf{F1} \\
\midrule
0.50 (default) & 0.218 & 0.875 & 0.313 \\
0.70 & 0.412 & 0.723 & 0.498 \\
0.78 (optimal) & 0.561 & 0.608 & 0.582 \\
0.85 & 0.645 & 0.391 & 0.472 \\
0.90 & 0.712 & 0.198 & 0.295 \\
\bottomrule
\end{tabular}
\caption{CosineSim threshold sweep. Even after optimization, CosineSim remains the weakest baseline, confirming stagnation $\neq$ surface repetition.}
\label{tab:cosinesim}
\end{table}

%==============================================================================
% 5.11 DOWNSTREAM APPLICATION
%==============================================================================
\subsection{Early Termination Simulation}

To assess downstream utility, we simulate probe-guided early termination of agent trajectories (Table~\ref{tab:early-term}).

\begin{table}[t]
\centering
\small
\begin{tabular}{lccc}
\toprule
\textbf{Threshold} & \textbf{Success Rate} & \textbf{Tokens Saved} & \textbf{False Term.} \\
\midrule
Baseline (no term.) & 28.5\% & 0\% & 0\% \\
Probe $t{=}$0.8 & 26.8\% & 18.5\% & 6.2\% \\
Probe $t{=}$0.6 & 24.2\% & 32.1\% & 14.8\% \\
Probe $t{=}$0.4 & 19.8\% & 48.7\% & 28.5\% \\
Oracle & 28.0\% & 35.2\% & 0.5\% \\
\bottomrule
\end{tabular}
\caption{Early termination. Conservative threshold ($t{=}$0.8) saves 18.5\% tokens with minimal success-rate impact ($-$1.7pp). The oracle gap highlights the cost of imperfect detection.}
\label{tab:early-term}
\end{table}

A conservative threshold ($t{=}$0.8) saves 18.5\% of tokens with only 1.7pp success-rate loss and 6.2\% false termination rate.
However, aggressive thresholds lead to unacceptable false termination rates (28.5\% at $t{=}$0.4).
The oracle shows that a perfect detector could save 35.2\% of tokens with negligible cost, illustrating the theoretical value of stagnation detection---and the practical gap that current probe accuracy leaves.

%==============================================================================
% 5.12 INTERPRETABILITY
%==============================================================================
\subsection{Interpretability Analysis}

t-SNE visualization of GRU hidden states shows moderate separation between progress and stagnation clusters (silhouette score 0.31), with substantial overlap at class boundaries.
Attention weight analysis reveals no clear ``stagnation onset detection'' pattern: weights distribute relatively uniformly across the window, with a slight bias toward the middle steps rather than concentrating on transition points.
The GRU appears to perform coarse-grained pattern matching rather than fine-grained temporal reasoning, consistent with its moderate gold-set performance.

%==============================================================================
% 6. DISCUSSION
%==============================================================================
\section{Discussion}

\paragraph{What data programming gets wrong on sequential data.}
Our findings reveal a fundamental tension: LFs on agent trajectories exhibit systematic temporal dependencies (leading/lagging structure) that violate the i.i.d.\ assumption, yet explicitly modeling these dependencies (via temporal label model) yields no improvement over correlation-aware methods (FlyingSquid) that are agnostic to temporal order.
The explanation is that within $k{=}5$ windows, temporal dynamics are averaged out---the signal exists at a finer granularity than the task formulation captures.

\paragraph{The leaky-heuristic trap.}
Our strongest LF (\texttt{outcome\_bp}) is also the leakiest, contributing 7.4pp while encoding trajectory-level outcome rather than window-level progress.
This creates a pernicious feedback loop: the most informative heuristic inflates both label quality and downstream model performance, making the pipeline appear more effective than it is.
We estimate ``clean'' probe performance (without \texttt{outcome\_bp} influence) at 0.68--0.70 on the gold set---near the structural baseline.

\paragraph{When probes are (not) useful.}
The probe's gold-set F1 (0.72) is clearly below LLM-as-judge CoT (0.855).
Its value proposition is limited to scenarios where LLM evaluation is economically infeasible: continuous monitoring of $>$100K windows, where the 615$\times$ cost advantage dominates.
For smaller-scale evaluation, LLM-as-judge is strictly preferable.

\paragraph{Inherent task ambiguity.}
Human annotator agreement ($\kappa{=}0.52$) sets a ceiling on how cleanly any method can separate progress from stagnation.
The boundary is genuinely fuzzy---exploration can look like stagnation, and partial progress resists binary classification.
Future work should consider ordinal or continuous progress scales rather than binary labels.

%==============================================================================
% 7. CONCLUSION
%==============================================================================
\section{Conclusion}

We present an empirical study of weak supervision for progress stagnation detection in LLM agent trajectories.
Our original hypotheses---that temporal label modeling and neural probes would advance the state of detection---do not hold.
Instead, we contribute:
(1)~a systematic characterization of data programming failure modes on sequential agent data, including i.i.d.\ violation through LF temporal dynamics and outcome leakage through backward-propagating heuristics;
(2)~a stagnation taxonomy and window-level formulation that provide a concrete task definition;
(3)~empirical evidence that GRU temporal modeling provides disproportionate robustness to label noise (+6.9pp over mean-pooling at high noise).
These findings provide cautionary lessons and empirical foundations for future work on agent monitoring via weak supervision.

\section*{Limitations}

Our study has several limitations.
(1)~All experiments use software engineering agent trajectories (SWE-bench); generalization to other domains (web navigation, embodied agents) is unknown.
(2)~Cross-agent transfer is evaluated on only two implementations (SWE-agent and OpenHands).
(3)~The gold set (420 windows) is small relative to the full dataset; larger annotation efforts could yield tighter performance estimates.
(4)~Human annotator agreement ($\kappa{=}0.52$) questions the reliability of the gold set itself as ground truth.
(5)~We do not evaluate online deployment or integration with meta-controllers.

\bibliographystyle{plainnat}
\begin{thebibliography}{20}

\bibitem[Yang et~al.(2024)]{yang2024swe}
John Yang, Carlos~E. Jimenez, Alexander Wettig, Kilian Lieret, Shunyu Yao, Karthik Narasimhan, and Ofir Press.
\newblock {SWE}-agent: Agent-computer interfaces enable automated software engineering.
\newblock \emph{arXiv preprint arXiv:2405.15793}, 2024.

\bibitem[Wang et~al.(2024)]{wang2024opendevin}
Xingyao Wang, Boxuan Li, Yufan Song, Frank~F. Xu, Xiangru Tang, Mingchen Zhuge, Jiayi Pan, Yueqi Song, Bowen Li, Jaskirat Singh, et~al.
\newblock {OpenDevin}: An open platform for {AI} software developers as generalist agents.
\newblock \emph{arXiv preprint arXiv:2407.16741}, 2024.

\bibitem[Ratner et~al.(2017)]{ratner2017snorkel}
Alexander Ratner, Stephen~H. Bach, Henry Ehrenberg, Jason Fries, Sen Wu, and Christopher R{\'e}.
\newblock Snorkel: Rapid training data creation with weak supervision.
\newblock In \emph{Proc.\ VLDB}, 2017.

\bibitem[Fu et~al.(2020)]{fu2020fast}
Daniel~Y. Fu, Mayee~F. Chen, Frederic Sala, Sarah~M. Hooper, Kayvon Fatahalian, and Christopher R{\'e}.
\newblock Fast and three-rious: Speeding up weak supervision with triplet methods.
\newblock In \emph{Proc.\ ICML}, 2020.

\bibitem[Shin et~al.(2022)]{shin2022universalizing}
Jieyu Shin, Gilles Baechler, Weilin Xu, Ranjay Krishna, and Christopher R{\'e}.
\newblock Universalizing weak supervision.
\newblock In \emph{Proc.\ ICLR}, 2022.

\bibitem[Lightman et~al.(2023)]{lightman2023lets}
Hunter Lightman, Vineet Kosaraju, Yura Burda, Harri Edwards, Bowen Baker, Teddy Lee, Jan Leike, John Schulman, Ilya Sutskever, and Karl Cobbe.
\newblock Let's verify step by step.
\newblock In \emph{Proc.\ ICLR}, 2023.

\bibitem[Wang et~al.(2024b)]{wang2024math}
Peiyi Wang, Lei Li, Zhihong Shao, R.X. Xu, Damai Dai, Yifei Li, Deli Chen, Y.\ Wu, and Zhifang Sui.
\newblock Math-{S}hepherd: Verify and reinforce {LLM}s step-by-step without human annotations.
\newblock In \emph{Proc.\ ACL}, 2024.

\bibitem[Shinn et~al.(2023)]{shinn2023reflexion}
Noah Shinn, Federico Cassano, Ashwin Gopinath, Karthik Narasimhan, and Shunyu Yao.
\newblock Reflexion: Language agents with verbal reinforcement learning.
\newblock In \emph{Proc.\ NeurIPS}, 2023.

\bibitem[Yao et~al.(2024)]{yao2024tree}
Shunyu Yao, Dian Yu, Jeffrey Zhao, Izhak Shafran, Tom Griffiths, Yuan Cao, and Karthik Narasimhan.
\newblock Tree of thoughts: Deliberate problem solving with large language models.
\newblock In \emph{Proc.\ NeurIPS}, 2024.

\bibitem[Le et~al.(2022)]{le2022coderl}
Hung Le, Yue Wang, Akhilesh~Deepak Gotmare, Silvio Savarese, and Steven~Chu-Hong Hoi.
\newblock {CodeRL}: Mastering code generation through pretrained models and deep reinforcement learning.
\newblock In \emph{Proc.\ NeurIPS}, 2022.

\bibitem[Zheng et~al.(2024)]{zheng2024judging}
Lianmin Zheng, Wei-Lin Chiang, Ying Sheng, Siyuan Zhuang, Zhanghao Wu, Yonghao Zhuang, Zi~Lin, Zhuohan Li, Dacheng Li, Eric~P. Xing, et~al.
\newblock Judging {LLM}-as-a-judge with {MT-Bench} and {Chatbot Arena}.
\newblock In \emph{Proc.\ NeurIPS}, 2024.

\bibitem[CDCS(2024)]{cycles2024}
{Agent Cycle Detection Consortium}.
\newblock Detecting and mitigating repetitive cycles in {LLM} agent trajectories.
\newblock \emph{arXiv preprint}, 2024.

\bibitem[SpecRA(2024)]{specra2024}
{SpecRA Team}.
\newblock Spectral analysis for anomaly detection in agent trajectories.
\newblock \emph{arXiv preprint}, 2024.

\end{thebibliography}

\end{document}
